\subsection{Three focused questions about social context}
Gini measures income inequality, with larger values indicating a more unequal distribution. The Gender Inequality Index (GII) measures gender disadvantage in health, empowerment and employment. PPP income adjusts purchasing power for differences in the cost of living. Homicide is measured as deaths per 100,000 people. None of these indicators directly measures what people believe about a song.
Songs can describe, criticize or imagine social conditions. Their chart presence also depends on which genres and international hits circulate. Research on cross-country listening links musical preferences to language, culture and social networks~\cite{liu}. Work on lyrical change shows why the supply and composition of songs matter~\cite{varnum,parada}. These findings motivate a selection-based interpretation: a country indicator and a lyric theme can be associated because markets contain different repertoires, even when the lyrics do not mirror residents' lives.

We therefore narrow the contextual analysis to three questions. First, does the relationship between materialism and income inequality remain after accounting for economic resources and repertoire composition? Wealth references may express aspiration or status display, but they can also be concentrated in commercially successful genres. Second, does the relationship between misogyny and the Gender Inequality Index (GII) survive adjustment for the same differences? Depicting a gender stereotype and endorsing it are separate processes, so national gender conditions need not produce more hostile lyrics. Third, does the relationship between lyrical violence and homicide persist after adjustment? Violence can be narrated, condemned or dramatized, while a market can import songs from places with very different conditions. These are competing explanations to examine, rather than reasons to predict a particular sign.

The three corresponding hypotheses concern an association that remains after accounting for measured selection: H1 pairs materialism with inequality, H2 pairs misogyny with gender conditions, and H3 pairs lyrical violence with homicide. A direct social-reflection account would predict a remaining contextual association. A repertoire-selection account would predict attenuation once economic resources, genre and shared circulation are considered. The sign is left open because depiction, criticism and aspiration can operate in different directions. These hypotheses were formulated after inspecting the original descriptive results. The analysis is exploratory and has not been preregistered. The complete correlation matrix remains a supplementary descriptive resource, rather than the basis for selecting additional claims.

To keep the music source constant, the focused models use only the 2021 MGD Spotify Top-200 repertoire. Every market receives equal weight. Each outcome is the percentage of classified recordings with clearly present materialism, misogyny or violence. The focal predictors are Gini, GII and log(1 + homicide rate), respectively. We compare an unadjusted model with a model that also includes log PPP income, internet access, rap/hip-hop share, the share without genre tags, and the share of recordings charting in at least five panel markets. The last variable measures shared circulation, not foreign artist origin. Rap share addresses one important genre contrast but cannot replace complete genre composition. Actual lyric language is unavailable and is not inferred from a country's dominant language.

Within each question, both models use exactly the same markets with complete data. Coefficients describe a change in thematic share, in percentage points, per one standard deviation of the contextual predictor in that sample. We use linear regression with HC3 standard errors, which allow residual variation to differ across markets, and Student-$t$ intervals. Holm correction covers the three adjusted tests. Shared songs can still produce residual dependence between markets, and missing language and genre metadata leave confounding unresolved. The models examine whether the original associations are sensitive to measured composition. They do not identify causal effects or validate prediction outside the sample.

\begin{table*}[t]
\caption{Focused social associations in the 2021 MGD Top-200 repertoire. Each entry is a coefficient with its 95\% HC3 interval, in percentage points per sample standard deviation of context. Unadjusted and adjusted estimates use identical markets within each row. Adjusted models include income, internet access and the three repertoire controls described in the text.}
\label{tab:context}
\centering\small
\begin{tabularx}{\textwidth}{@{}Xrlll@{}}
\toprule Theme / context & Markets & Unadjusted [95\% interval] & Adjusted [95\% interval] & Holm $p$ \\
\midrule
@@FOCUSED_SOCIAL_ROWS@@
\bottomrule
\end{tabularx}
\end{table*}

The three negative unadjusted associations become small positive estimates after adjustment, and all adjusted intervals include zero (Table~\ref{tab:context}). Materialism--Gini changes from $-2.90$ to $+0.43$ percentage points per standard deviation. Misogyny--GII changes from $-0.94$ to $+0.47$, and violence--homicide from $-2.09$ to $+1.19$. This sensitivity shows why the original negative correlations cannot be read as evidence that inequality, gender disadvantage or homicide protect against the corresponding lyrical themes. It also does not establish a positive relationship: the adjusted estimates are uncertain, and the controls are incomplete.

Brazil's same-year contextual values provide a concrete reference: Gini is 52.9, GII is 0.405, internet access is 80.7\%, and homicide is 21.7 per 100,000 people. These indicators describe social conditions. The lyrical percentages describe selected chart repertoires. Keeping the two units distinct makes the comparison useful without treating songs as a survey of the population.
